%%%%%%%%%%%%%%%%%%%%%%%%%%%%%%%%%%%%%%%%%%%%%%%%%%%%%%%%%%%%%%%%%%%%%%%%%%%%%%%
% Theorem stuff, use "thm" for italic fonts inside theorem environments and   %
% "def" for normal fonts. These differ from the originals: newline is created %
% after the complete name (including optional parenthesis) of the theorem.    %
%%%%%%%%%%%%%%%%%%%%%%%%%%%%%%%%%%%%%%%%%%%%%%%%%%%%%%%%%%%%%%%%%%%%%%%%%%%%%%%

\newtheoremstyle{thm}   % name
{15pt}                  % Space above
{10pt}                  % Space below
{\itshape}              % Body font
{}                      % indent amount
{\bfseries}             % Theorem head font
{}                      % Punctuation after theorem head
{\newline}              % space after theorem head
{}                      % Theorem head spec

\newtheoremstyle{def}   % name
{15pt}                  % Space above
{10pt}                  % Space below
{}                      % Body font
{}                      % indent amount
{\bfseries}             % Theorem head font
{}                      % Punctuation after theorem head
{\newline}              % space after theorem head
{}                      % Theorem head spec

%%%%%%%%%%%%%%%%%%%%%%%%%%%%%%%%%%%%%%%%%%%%%%%%%%%%%%%%%%%%%%%%%%%%%%%%%%%%%%%
% Forces use of "Number.Subnumber Name" (Wrobel style ;) ). Enumeration       %
% follows the current chapter NOT the current section.                        %
% Do NOT use this, if you don't use chapters. Swap the comments below.        %
%%%%%%%%%%%%%%%%%%%%%%%%%%%%%%%%%%%%%%%%%%%%%%%%%%%%%%%%%%%%%%%%%%%%%%%%%%%%%%%

\swapnumbers
\theoremstyle{thm}
\newtheorem{sat}{Satz}[chapter]

\theoremstyle{def}
\newtheorem{bem}{Bemerkung}[section]

\theoremstyle{def}
\newtheorem{dfn}[sat]{Definition}

\theoremstyle{thm}
\newtheorem{kor}[sat]{Korollar}

\theoremstyle{thm}
\newtheorem{lem}[sat]{Lemma}

\theoremstyle{thm}
\newtheorem{defsat}[sat]{Definition/Satz}

\theoremstyle{thm}
\newtheorem{beisp}[sat]{Beispiel}

% These commands (definition through korollar) can be used to ensure the
% following:
% - you open and close a new "something", name it (#1) and label it (#2).
% - the label is set in the right place (i. e. hyperref will jump to the right
%   position; the label is "#2"
% - a tableofcontents-entry is created, which respects the same enumeration

\newenvironment{definition}[2] % 1=Name, 2=Label
{\begin{dfn}[#1]
\label{#2}
\addcontentsline{toc}{section}{\protect\numberline{\ref{#2}} {#1}} }
{\end{dfn}}

\newenvironment{lemma}[2] % 1=Name, 2=Label
{\begin{lem}[#1]
\label{#2}
\addcontentsline{toc}{section}{\protect\numberline{\ref{#2}} #1}}
{\end{lem}}

\newenvironment{satz}[2] % 1=Name, 2=Label
{\begin{sat}[#1]
\label{#2}
\addcontentsline{toc}{section}{\protect\numberline{\ref{#2}} #1}}
{\end{sat}}

\newenvironment{bemerkung}[2] % 1=Name, 2=Label
{\begin{bem}[#1]
\label{#2}
\addcontentsline{toc}{section}{\protect\numberline{\ref{#2}} #1}}
{\end{bem}}

\newenvironment{korollar}[2] % 1=Name, 2=Label
{\begin{kor}[#1]
\label{#2}
\addcontentsline{toc}{section}{\protect\numberline{\ref{#2}} #1}}
{\end{kor}}

\newenvironment{defsatz}[2] % 1=Name, 2=Label
{\begin{defsat}[#1]
\label{#2}
\addcontentsline{toc}{section}{\protect\numberline{\ref{#2}} #1}}
{\end{defsat}}

\newenvironment{beispiel}[2] % 1=Name, 2=Label
{\begin{beisp}[#1]
\label{#2}
\addcontentsline{toc}{section}{\protect\numberline{\ref{#2}} #1}}
{\end{beisp}}